\documentclass[track=junior,level=3]{kidscode}
%% \gl{term}{week you first meet it}{meaning}
%% Built from the vocabulary boxes of the workbook chapters (one entry per word).
\newcommand{\gl}[3]{\par\noindent\needspace{3\baselineskip}{\bfseries #1}\hfill{\footnotesize\color{kcgray}#2}\par\nopagebreak
  \noindent #3\par\vspace{6pt}}
\newcommand{\letter}[1]{\par\needspace{6\baselineskip}\vspace{4pt}\noindent{\Large\bfseries\color{kcaccent}#1}\par\nopagebreak\vspace{2pt}}
\begin{document}
\kctitle{Word Book}{Every new word in Bridge to Python, from A to Z}

\begin{grownup}[How to use the word book]
Read a word and its meaning aloud, then ask your child to show you the word on the computer or act
it out. The gray label says which week the word first comes up (W1 = week 1).
\end{grownup}


\letter{B}
\gl{bar graph}{W3}{A picture of data. The tallest bar has the most.}

\letter{D}
\gl{data}{W3}{Information we collect, like answers to a question.}

\letter{E}
\gl{error message}{W6}{Red text from Python. It says which line has a problem.}
\gl{example}{W4}{One thing we show the computer, marked YES or NO.}

\letter{F}
\gl{feature}{W4}{Something the computer looks at, like legs or wings.}

\letter{G}
\gl{graduate}{W10}{Someone who has finished a level and is ready for the next one.}

\letter{I}
\gl{indent}{W7}{Spaces at the start of a line. They put the line inside the loop.}

\letter{L}
\gl{loop}{W7}{Code that repeats. In Python, \texttt{for} makes a loop.}

\letter{N}
\gl{nested repeat}{W1}{A repeat inside another repeat.}

\letter{P}
\gl{part}{W2}{A small piece of a big program. You can build and test it on its own.}
\gl{partner check}{W9}{A friend tests your project against your plan.}
\gl{pattern}{W1}{Something that happens again and again in the same way.}
\gl{pen size}{W8}{How thick the turtle's line is. Bigger numbers draw thicker lines.}
\gl{plan card}{W9}{A plan on paper: what you will build, in parts.}
\gl{prediction}{W4}{The computer's guess about something new.}
\gl{Python}{W5}{A coding language that uses words instead of blocks.}

\letter{R}
\gl{remix}{W8}{A new version of someone's program. Say whose idea it was!}

\letter{S}
\gl{showcase}{W10}{A little show where we share what we made.}
\gl{step list}{W2}{The steps for one part, written in order.}

\letter{T}
\gl{tally mark}{W3}{One line for each answer. Five marks look like a gate.}
\gl{test}{W2}{Run it and check that it does what the plan says.}
\gl{Thonny}{W6}{The app where we write and run Python.}
\gl{typo}{W6}{A small spelling mistake, like \texttt{forwad}.}

\letter{W}
\gl{word twin}{W5}{A Python line that does the same job as a Scratch block.}

\clearpage
\section{My own words}
Draw or write a new word, and what it means to you.
\begin{center}\renewcommand{\arraystretch}{3}
\begin{tabularx}{\linewidth}{|p{4.5cm}|X|}
\hline
\textbf{Word} & \textbf{What it means (draw or write)}\\ \hline
 & \\ \hline
 & \\ \hline
 & \\ \hline
 & \\ \hline
 & \\ \hline
\end{tabularx}
\end{center}
\end{document}
